\section{Related Work}
\label{related_work}

\subsection{Traditional Requirement Elicitation}
Traditional requirement elicitation relies on interviews, workshops, scenarios, questionnaires, and document analysis to externalize stakeholder needs.
These techniques provide rich human context but depend heavily on analyst expertise and may miss implicit domain assumptions.

\subsection{Requirement Clarification}
Requirement clarification addresses incomplete, vague, or ambiguous stakeholder statements through follow-up questions and iterative refinement.
This paper treats clarification as a focused elicitation activity rather than a downstream specification or implementation task.

\subsection{LLM-Based Requirement Elicitation}
LLMs can support elicitation by summarizing stakeholder input, proposing requirement candidates, and generating follow-up questions.
The key limitation is that generated questions may be generic, weakly connected to missing information, or unsupported by explicit requirement knowledge.

\subsection{Multi-Agent Collaboration}
Multi-agent collaboration assigns elicitation responsibilities to specialized roles.
In this paper, collaboration is used narrowly for requirement clarification: one agent conducts the interview, one simulates or represents end-user perspectives, and one analyzes requirement quality during the dialogue.

\subsection{Human-in-the-Loop Requirement Engineering}
Human-in-the-loop requirement engineering keeps stakeholders and analysts responsible for validating questions, answers, and clarified requirement statements.
This line of work motivates interactive clarification workflows in which AI agents assist elicitation without replacing stakeholder judgment.
